\documentclass[11pt]{article}
\usepackage[margin=1in]{geometry}
\usepackage{lipsum,tikz,xcolor}
\usepackage[colorlinks=true,linkcolor=blue,urlcolor=teal,bookmarks=true,bookmarksopen=true]{hyperref}
\title{NightPdf Sample Document}
\author{Test Fixture}
\begin{document}
\maketitle
\tableofcontents
\noindent Jump to the figure in Section~\ref{sec:fig} on page~\pageref{fig:circles}, or read the conclusion in Section~\ref{sec:end}.
See also \href{https://www.apple.com}{an external link}.
\newpage
\section{Introduction}\label{sec:intro}
\lipsum[1-4]\footnote{A footnote that links back to Section~\ref{sec:end}.}
\subsection{Background}
\lipsum[5-9]
\subsection{Motivation}
\lipsum[10-14]
\section{Method}
\lipsum[15-20]
\subsection{Details}
\lipsum[21-26]
\section{Figures}\label{sec:fig}
\begin{figure}[h]\centering
\begin{tikzpicture}
\fill[red!80] (0,0) circle (1.2); \fill[green!60!black] (2.8,0) circle (1.2); \fill[blue!80] (5.6,0) circle (1.2);
\fill[orange] (8.4,0) circle (1.2);
\end{tikzpicture}
\caption{Colored circles to check dark-mode rendering.}\label{fig:circles}
\end{figure}
\lipsum[27-34]
\section{Conclusion}\label{sec:end}
Go back to the \hyperref[sec:intro]{introduction} or the figure on page~\pageref{fig:circles}.
\lipsum[35-45]
\end{document}
